\documentclass[11pt,a4paper,oneside,openright]{report}
%\documentclass[11pt,a4paper,openright]{report}
\usepackage{anysize}
\usepackage{vpcoestyle}   %% Thesis Title Page
\usepackage{latexsym}
%\usepackage{pictex}
\usepackage{fancyvrb}
\usepackage{fancybox}
\usepackage{fancyhdr}
\usepackage{color}
\usepackage{setspace}
\usepackage{amssymb}
\usepackage{multicol}
\usepackage{graphicx}
\usepackage{subfigure}
\usepackage{rawfonts}
\usepackage{setspace}

%\usepackage{newcent} %the options are:  palatino, times, newcent, bookman

\marginsize{1.5in}{1in}{0.75in}{0.75in}

\usepackage[pdftex,
        colorlinks=true,
        urlcolor=black,       % \href{...}{...} external (URL)
        filecolor=black,     % \href{...} local file
        linkcolor=black,       % \ref{...} and \pageref{...}
	citecolor=black,
        pdftitle={Seminar Synopsis},
        pdfauthor={Pratyay Dhond},
        pdfsubject={Seminar},
        pdfkeywords={Flutter},
%        pagebackref,
%        pdfpagemode=true,
        bookmarksopen=true]{hyperref}
\pdfcompresslevel=9


\thisfancypage{%
  \setlength{\fboxsep}{14pt}\shadowbox}{}

\title{Native Application Development using flutter}
\author{Mast. Pratyay Dhond}
\date{OCT 26, 2021}

\begin{document}
\def\what{SYNOPSIS}
\def\title{Native Application Development using flutter}

\def\degree{{\bf \large Diploma in Engineering\\(Information Technology)}}
\def\first{\large  Mast. Pratyay Prasad Dhond }
\def\one{\Large 1907011}

\def\guide{\Large Dr. A R Mahajan}
\def\coguide{Prof.XYZ }
\pagenumbering{roman}

%\ifpdf
\pdfbookmark[1]{Title}{title} %%This is pdf specefic
%\fi

%--------------Title Page-------------
\vpcoetitlepage
\cleardoublepage
\thispagestyle{empty}
\thisfancypage{%
  \setlength{\fboxsep}{14pt}\shadowbox}{}
\begin{center} \label{chap:approval}
{\bf \huge {Approval Sheet}} \\

\vspace{0.8in}
\end{center}
\renewcommand{\baselinestretch}{1.5}
\hspace{0.2in} \large {The synopsis entitled {\bf Native Application Development using Flutter} submitted by \par{\hspace{0.4in}{\bf Mast. Pratyay Prasad Dhond}. {\bf Roll No:1907011}\par{\hspace{0.4in}\\ is approved for the partial fulfilment of the requirement for the award of {\bf Diploma} in {\bf Information Technology}}\\

\vspace{1.5in}
\begin{multicols}{2}
\begin {center}
 {\bf Dr. A R Mahajan}\\ Guide\\%\\Department of Computer Engineering\\.\\
 {\bf Prof.L.D.Vilhekar}\\Project Coordinator\\%\\Department of Computer Engineering\\.\\
\end {center}
\end{multicols}
\vskip 0.6in
\begin{center}
\begin{multicols}{2}

.\\{\bf Seal/Stamp}\\
{\bf \large Dr. A.R.Mahajan}\\{\bf HOD}\\
\end{multicols}
\end{center}


\newpage

%

%%---------------- A B S T R A C T -----------------
\onehalfspacing
\cleardoublepage
\vspace*{2in}
\section*{\begin{center} {\Huge{Abstract}}\end{center}}
%\thispagestyle{empty}
\addcontentsline{toc}{chapter}{\quad\,\,{Abstract}}

 \hspace{0.2in}Flutter is a cross-platform User Interface development framework that is used to develop cross platform applications using a single code base. Flutter can be used to build beautiful, natively compiled application for various platforms such as Android, iOS, Windows, Mac, Linux, Web, etc. using a single codebase. 

 \hspace{0.2in}The key features of the Flutter Framework are Fast Development using the Hot Reload, Expressive and Flexible User Interface(UI) and the Native Performance.

\hspace{0.2in}Fast Development : Flutter's hot reload helps you quickly and easily experiment, build UIs, add features, and fix bugs faster. Experience sub-second reload times without losing state on emulators, simulators, and hardware.

\hspace{0.2in} Expressive, beautiful UIs: Flutter comes with easy to build, highly customizable UI widgets which can be used for easy development. Flutter also supports custom widget building using other widgets.

\hspace{0.2in} Native Performance : Flutter’s widgets incorporate all critical platform differences such as scrolling, navigation, icons and fonts to provide full native performance on both iOS and Android. 

%-----------------Table of Contents-------------------
%\renewcommand{\tableofcontents}{Table of Contents}
\newpage
%\ifpdf
\pdfbookmark[1]{Table of Contents}{contents} %%This is pdf specefic
%\fi
\tableofcontents
%\thispagestyle{empty}
\addtocontents{toc}{\protect\thispagestyle{empty}}
\newpage

\listoffigures
 2.3.4.1 Dash: Official Mascot of Dart and Flutter
 
%\thispagestyle{empty}
\addcontentsline{toc}{chapter}{\quad\,\,{List of Figures}}
\addtocontents{lof}{\protect\thispagestyle{empty}}
%\newpage




%----------------------- Chapter 1 ---------------------
\chapter{Introduction}
\label{chap:intro}
\pagenumbering{arabic}

\section{Introduction}
\hspace{0.2in} If you ask ten different mobile developers how they develop their mobile applications for Android or iOS devices, you'll probably get 10 different answers. This use of different languages for different Operating systems creates the need of employment of development teams for each operating system like Android, iOS, Windows, Mac, Linux, as well as web-apps,etc. This severely affects the cost of the project to build. 
				
\hspace{0.2in} To avoid this reason there was a need for cross platform development from the old-fashioned porting of apps. For solving this problem, Google developed the flutter framework to make cross platform development easier, with the view of -'One codebase, multiple platforms!'.
			

\begin{itemize}

\item {\bf History of Flutter : }
\end{itemize}
\hspace{0.2in} Flutter began it's life in the year 2017 under the name 'sky' at the dart developer summit.At first, it ran only on google's own own Android operating system, but before long was ported to Apple's ios. 

\hspace{0.2in} Various preveiw versions of Flutter were released subsequent to its initial announcement, culminating in the December  4th,2018, release of Flutter 1.0, the first "Stable" release. One of the main goals of Flutter as stated was, 'being able to render app UIs at consistent 120fps no matter what."

\hspace{0.2in} Fltter iffers twi sets if wudgets : Material design widgets and Cupertino design widgets.

\section{Flutter Keywords}
 \begin{enumerate}
 \item UI : User Interface
 \end{enumerate}

%\cleardoublepage

\chapter{Problem Statement}\label{chap:project outline}

\section{Problem Statement}
 Native Application Development using Flutter.
\section{Solving Approach}
	Flutter is a comparatively new framework in the cross platform development field, released recently in the December of 2018. From it's release it has received a lot of adoration from the developer community, as it is build on Dart a language developed by Google itself.  

\section{Outcomes}
\hspace{0.2in} These are the different tags which will be included in the Seminar :
\begin{enumerate}
\item Home :
\begin{itemize}
\item Flutter
\item Dart
\item Native App
\item OS
\item Web
\item Cross Platform
\item Contents\\
\\
\\
\end{itemize}
\item History :  
%This will contain the details of Flutter such as its history and the timeline of flutter.
	\begin{itemize}
	\item The first version of flutter was known by the codename as 'sky' and ran on the Android OS. It was unvailed at the Dart developer summit in the year 2015.
	\item The main intent of flutter SDK was to render at a consistent 150 frames per second.
	\item On 4th December 2018, the first stable version of Flutter , v1.0 was released at the flutter live event.
	\item On September 8th 2021, Dart SDK version 2.14 and flutter SDK version 2.5 were released by Google.
	\end{itemize}

\item Importance :
\begin{itemize}
\item Flutter framework is very important and is trending for a reason: it's an efficient way of building great-performing and beautiful apps.
\item Flutter allows you to create cross-platform apps that provide native performance. 
\item Apps created with Flutter feature beautiful and intuitive design and are able to run animations smoothly. 
\item   It comes with a wide variety of dependencies which are easily accessible and well documented. 
\item The dependencies can be found on 'https://pub.dev/'.
\end{itemize} 

\item Photo Gallery :\\
\hspace{0.2in}
\includegraphics[scale=0.15]{dash.png}\\
\hfill \includegraphics[scale=0.8]{flutter.png} \hfill\\ 
\hfill \includegraphics[scale=0.14]{dart.png} \hfill\\



\item News :\\
\hspace{0.2in}It contains the information about the latest updates, news and features added in dart as well as flutter.

\item Overview:\\
\hspace{0.2in} This part gives a short overview of the whole presentation.

\item My opinion :\\
\hspace{0.2in} This part contains my personal opinion about the future of Dart and Flutter.

\end{enumerate}





\chapter{References}
\begin{enumerate}

\item https://flutter.dev/\\
{\em "This is the official website of Flutter."}

\item https://pub.dev/\\
{\em "This is a website where developers can access various open source dependencies and implement them in their project."}

\item https://dart.dev/\\
{\em "This is the official website for Dart programming language."}

\item Practical Flutter - Frank Zammetti\\
{\em "This is a great book for anyone who wants to learn to code using flutter framework and build cross platform apps."}

\end{enumerate}










\end{document}
 